\newglossaryentry{arquitecturamedallon}{
	name={Arquitectura medallón},
	description={Patrón de diseño de \textit{pipelines} de datos que organiza el flujo de información en capas sucesivas de refinamiento (Bronce, Plata y Oro), desde el dato crudo hasta las vistas listas para consumo}
}

\newglossaryentry{calibracion}{
	name={Calibración de probabilidades},
	description={Procedimiento que ajusta las salidas de un modelo de clasificación para que reflejen probabilidades reales, permitiendo aplicar umbrales de confianza sobre sus predicciones}
}

\newglossaryentry{chatbot}{
	name={Chatbot},
	description={Interfaz de software que permite interactuar con un sistema mediante una conversación en lenguaje natural, en lugar de formularios o menús tradicionales}
}

\newglossaryentry{crispdm}{
	name={CRISP-DM},
	description={\textit{Cross Industry Standard Process for Data Mining}: modelo de proceso estándar para proyectos de analítica y minería de datos, estructurado en seis fases iterativas}
}

\newglossaryentry{datomaestro}{
	name={Dato maestro},
	description={Información que describe una entidad central de negocio (como un material, cliente o proveedor) que es referenciada de forma persistente por múltiples procesos y sistemas}
}

\newglossaryentry{duplicadodifuso}{
	name={Duplicado difuso},
	description={Par de registros que refieren al mismo objeto pero cuya descripción textual no es idéntica, por diferencias de abreviatura, orden de palabras o separadores}
}

\newglossaryentry{erp}{
	name={ERP},
	description={\textit{Enterprise Resource Planning}: sistema de información que integra los procesos centrales de una organización (finanzas, inventario, compras, entre otros) en una plataforma única}
}

\newglossaryentry{f1macro}{
	name={F1 macro},
	description={Métrica de desempeño de un clasificador que promedia el F1 de cada clase por igual, sin ponderar por su frecuencia, lo que penaliza el bajo desempeño en clases minoritarias}
}

\newglossaryentry{gestor}{
	name={Gestor de datos maestros},
	description={Persona responsable de crear, validar y mantener los registros del catálogo de materiales, asegurando su unicidad, exactitud y correcta categorización}
}

\newglossaryentry{llm}{
	name={LLM},
	description={\textit{Large Language Model} o modelo de lenguaje de gran escala: modelo de aprendizaje automático entrenado sobre grandes volúmenes de texto, capaz de generar y transformar lenguaje natural}
}

\newglossaryentry{mdm}{
	name={MDM},
	description={\textit{Master Data Management} o gestión de datos maestros: disciplina encargada de mantener consistentes, correctas y libres de duplicidad las entidades centrales de una organización}
}

\newglossaryentry{pgtrgm}{
	name={pg\_trgm},
	description={Extensión de PostgreSQL que permite calcular la similitud entre cadenas de texto mediante la comparación de trigramas de caracteres}
}

\newglossaryentry{pipeline}{
	name={Pipeline de datos},
	description={Secuencia automatizada de procesos que extrae, transforma y carga datos desde una fuente de origen hasta su destino final de consumo}
}

\newglossaryentry{roi}{
	name={ROI},
	description={\textit{Return on Investment} o retorno de inversión: relación entre el beneficio obtenido de un proyecto y el costo invertido en él}
}

\newglossaryentry{svmlineal}{
	name={SVM lineal},
	description={Máquina de vectores de soporte con núcleo lineal: algoritmo de clasificación que busca el hiperplano que separa las clases maximizando el margen entre los puntos más cercanos de cada una}
}

\newglossaryentry{tfidf}{
	name={TF-IDF},
	description={\textit{Term Frequency – Inverse Document Frequency}: técnica que representa un texto de forma numérica ponderando cada término según su frecuencia relativa y su rareza en el corpus}
}

\newglossaryentry{topk}{
	name={Exactitud top-\textit{k}},
	description={Métrica que considera una predicción correcta si la clase real se encuentra entre las \textit{k} categorías de mayor probabilidad sugeridas por el modelo, no solo en la primera}
}

\newglossaryentry{trigrama}{
	name={Trigrama},
	description={Subsecuencia de tres caracteres consecutivos extraída de una cadena de texto, utilizada como unidad básica para calcular similitud textual}
}

\newglossaryentry{umbralconfianza}{
	name={Umbral de confianza},
	description={Valor mínimo de probabilidad que debe alcanzar una predicción del modelo para presentarse como sugerencia automática; por debajo de él, la decisión se deriva a revisión del gestor}
}

\newglossaryentry{unspsc}{
	name={UNSPSC},
	description={\textit{United Nations Standard Products and Services Code}: taxonomía abierta y multisectorial para la clasificación de productos y servicios, administrada por el Programa de las Naciones Unidas para el Desarrollo}
}

\newglossaryentry{vps}{
	name={VPS},
	description={\textit{Virtual Private Server} o servidor privado virtual: instancia de servidor alojada en la nube, utilizada para el despliegue de la plataforma}
}